\section{SAT and polynomial verifier reductions}\label{sec:complexity}

The root article states that satisfiability is NP-complete. Here that
statement receives a finite, constructive proof rather than being left as
an external complexity-theoretic assertion. The mechanism is the bounded
computation encoding behind Cook's theorem \cite{cook1971}, followed by the
explicit Tseitin construction proved earlier. We state the machine model,
the witnesses, and the uniformity needed for the iPRA reading.

\subsection{Representing the classes without quantifying over arbitrary sets}

Inputs are finite binary strings. A verifier presentation consists of a
fixed deterministic Turing machine $M$, a polynomial witness-length bound
$p(n)$, and a polynomial clock $T(n)$. Both polynomials have specified
natural coefficients and exponents; increase the bounds if necessary to
avoid negative or fractional coefficients. A witness has exactly $p(n)$
bits. A witness of length at most $q(n)$ may instead be represented by
$q(n)+1$ unary length-marker bits and $q(n)$ padded data bits; validate the
marker and ignore padding. This gives exact length $2q(n)+1$.
The accepted predicate is
\[
 L_{M,p,T}(x)\quad:\!\Longleftrightarrow\quad
 \exists w<2^{p(|x|)}\ \operatorname{Accept}_M(x,w,T(|x|))=1.
\]
The acceptance test is a primitive recursive bounded simulation. The
polynomially clocked verifier presentations are what ``NP'' means here.
A polynomial reduction is a specified deterministic string algorithm with
a specified polynomial running-time bound and a proved equivalence of
membership predicates. Thus the assertions below are uniform theorem
schemes about numerical programs and their certificates, not unformalized
quantification over a universe of arbitrary languages.

This verifier definition agrees with the finite branching nondeterministic
description: a nondeterministic run of polynomial length is specified by
its sequence of transition choices, each from a fixed finite list, and a
deterministic verifier replays those choices. Conversely, guess the
polynomial-length witness and run its deterministic verifier. Both
transformations have explicitly polynomial clocks. A multitape machine
may be replaced by a one-tape machine with polynomial overhead: record
the finitely many tracks and head markers in consecutive blocks, sweep
the at most polynomially many used cells to find and update the markers
for each simulated step, and repeat for the specified clock. The number
of tracks is fixed, so this uses at most a fixed polynomial number of
scanned cells. These are bounded simulations and therefore primitive
recursive constructions with bounded arithmetic invariants.

\begin{proposition}[Polynomial verifiability of SAT]\label{thm:sat-np}
Satisfiability of formulas, of Boolean circuits, and of clauses of width
at most three has a polynomially clocked verifier presentation. The
correctness of each presentation is provable in \ipra.
\end{proposition}
\begin{proof}
Parse the finite input and reject invalid syntax. Densely rename its
distinct atom indices to $0,\ldots,n-1$ by repeated comparisons; there are
at most as many as input symbol positions. For an input of bit length $m$,
a witness supplies exactly $m$ bits; consume its first $n$ bits and ignore
the remainder. Represent circuits as finite gate lists with all predecessor
indices smaller than their gate index, rejecting a malformed list.
Evaluate the formula or circuit in an order in which the children precede
their parent and accept precisely when the output is one. For clauses,
scan their literals and conjunction. The correctness equation is the
definition of evaluation and is proved by the node scan.
For a concrete bit-string model, repeated scans, comparisons of indices,
and copying a value from an earlier position each cost at most a polynomial
in the input length. A deliberately naive implementation with nested
linear scans is still polynomial; no constant-time random access to
unbounded integers is assumed. With repeated scans for parsing, comparisons
and lookup, a conservative one-tape clock $C(m+2)^8$ suffices for a fixed
implementation and fixed constant $C$: there are at most $m$ nodes and
indices of at most $m$ bits; using at most eight nested scans of length
$O(m+2)$ dominates parsing, renaming, table updates, and work-tape sweeps.
Witness length and total work therefore have polynomial bounds depending
only on input length.
\end{proof}

\subsection{Unrolling a verifier into a Boolean circuit}

Fix $M,p,T$ and an input $x$ of length $n$. Use a one-tape machine with a
finite tape alphabet $\Sigma$, a finite state set $Q$, and moves
$-1,0,+1$. Make the accepting and rejecting halting states stationary
self-loops, so acceptance by the clock is the same as acceptance at the
clock. The input consists of a fixed framing symbol, $x$, a separator,
the $p(n)$ witness bits, and another framing symbol; let its length be
$b=n+p(n)+3$. Put its first cell at offset $T(n)+1$ in an array of length
\[
                 \ell=b+2T(n)+3.
\]
There is enough blank padding that a head making at most $T(n)$ moves
cannot leave this array. The few extreme cells can be treated as fixed
blanks, since no head reaches them within the clock.

A cell state belongs to
\[
       K=\Sigma\ \sqcup\ (Q\times\Sigma).
\]
The second alternative records the unique head and its control state.
A legal configuration has exactly one tagged cell. There is a finite
local transition table $F:K^3\to K$ describing the next state of a cell
from its left neighbor, itself, and its right neighbor. To define it on a
legal triple, proceed as follows. If the central cell has the head,
consult $M$'s transition, write the specified symbol, and retain the head
tag only for a stationary move. If a neighboring head moves into this
cell, retain this cell's old tape symbol and attach the incoming control
state. If neither event occurs, retain the untagged symbol. Simultaneous
incompatible incoming heads cannot occur in a legal configuration;
on such triples and other invalid patterns, choose any fixed default.
This makes $F$ a total finite function. Its behavior on legal configurations
is exactly $M$'s next step.

Encode each cell by $k=|K|$ bits with exactly one bit true. For $d\in K$,
the output indicator for $d$ is the Boolean formula
\[
 \bigvee_{\substack{a,b,c\in K\\F(a,b,c)=d}}
                  (L_a\wedge C_b\wedge R_c).
\]
The input blocks are one-hot, so exactly one triple term is true and the
display gives precisely the indicated output. Empty disjunction means
false. Use binary conjunction and disjunction gates; each conjunction of
three bits needs two gates, and each disjunction of $r$ terms needs at
most $r$ gates with constants allowed. Across all $k$ output indicators,
a safe uniform bound is $4(k+1)^4$ gates per updated cell. The constants
depend only on the fixed machine.

\begin{theorem}[Bounded computation-to-circuit compiler]\label{thm:verifier-circuit}
There is a primitive recursive, polynomial-time compiler constructing a
circuit $C_x$ whose only free input bits are the witness bits, such that
\[
 \operatorname{ev}(w,C_x)=1
       \quad\Longleftrightarrow\quad
 \operatorname{Accept}_M(x,w,T(n))=1.
\]
For a fixed machine, a bound on its number of gates is
\[
 20(k+1)^4(T(n)+1)(\ell+1).
\]
The correctness equivalence, universally in the finite input and witness
codes, has an iPRA proof.
\end{theorem}
\begin{proof}
Construct the initial row of cell indicators from constants determined
by $x$, the witness inputs and their complements, and the unique initial
head tag. It is one-hot and encodes the prescribed initial configuration.
For each time $t<T(n)$ and each cell, append a fresh copy of the constant-size
local transition circuit just described, supplying constant blank neighbors
for any index just outside the array. Its input wires point only to
the preceding time row. Thus the circuit is acyclic and can be written in
time order. Finally take the disjunction of the indicators of accepting
head states in the last row, using a linear number of binary gates.

Induct on $t$ with the invariant that every cell's indicator block is
one-hot and agrees with the configuration after exactly $t$ simulated
steps. The base case is the initialization. In the step, all triples are
the triples of a legal configuration, so the finite table $F$ computes
the correct next cell, with exactly one output indicator true. The padding
bound ensures the boundary convention does not interfere. This proves
the invariant simultaneously for all cells. The final accept-disjunction
then has the claimed value. Every universal claim in this invariant is
bounded by $\ell,k$, and the time counter is bounded by $T(n)$, so it is
a bounded arithmetic formula and uses only permitted induction.

Initialization, the $T(n)\ell$ local copies, and the final disjunction fit
within the displayed generous gate bound. The output consists of gate
tags and predecessor indices of logarithmic bit length in that bound.
A sequential generator can write these indices with polynomially bounded
arithmetic and copying. Even representing its working tables by repeated
string scans introduces only a polynomial factor. Thus construction time
as well as output bit length is polynomial in $n$ for fixed $M,p,T$.
\end{proof}

\subsection{An actual NP-completeness proof}

The Tseitin encoding applies to acyclic circuits as well as syntax trees:
give each gate a fresh variable and impose the same local equivalence
clauses in dependency order. The proof that each gate variable must equal
its computed value uses exactly that order; sharing an input wire does
not affect it. Clauses for constants, negation, conjunction, disjunction
and implication have width at most three. Assert the output variable
by a unit clause. Retain an ordered list of all allocated witness and gate
variables, densely named $0,\ldots,v(\Phi_x)-1$; this defines
$v(\Phi_x)$. Thus unused allocated inputs also have assignment positions.
If one instead stores only occurring variables, fill any missing witness
positions with zero when projecting a satisfying assignment.

\begin{theorem}[Constructive NP-completeness of SAT and 3-SAT]
\label{thm:np-complete}
For each polynomially clocked verifier presentation $(M,p,T)$ there is a
polynomial-time primitive recursive map $x\mapsto\Phi_x$ into 3-CNF such
that \ipra\ proves
\[
 L_{M,p,T}(x)
       \quad\Longleftrightarrow\quad
 \exists a<2^{v(\Phi_x)}\ [\operatorname{ev}(a,\Phi_x)=1].
\]
Witness maps in both directions are explicit. Consequently SAT, circuit
SAT, and 3-SAT are NP-complete for these presentations under polynomial
many-one reductions.
\end{theorem}
\begin{proof}
Compile $C_x$ by Theorem~\ref{thm:verifier-circuit} and let $\Phi_x$ be
its Tseitin clauses with the output asserted. Given an accepting witness
$w$, evaluate the circuit in dependency order and extend $w$ by these
gate values. Every local equivalence and the output unit clause hold,
so this is a satisfying assignment of $\Phi_x$. Conversely, restrict any
satisfying assignment to the witness-input variables. Induction through
the gates, using the local equivalence clauses, forces every auxiliary
value to its actual circuit value. The asserted output is therefore the
value one, so the restricted witness is accepted by $M$. These two maps
are bounded table operations and establish the equivalence in iPRA.

If the circuit has $g$ gates, there are at most $3g+1$ local/output
clauses and a linear number of literal occurrences. Its binary encoding
has length $O((g+p(n)+1)\log(g+p(n)+2))$. The earlier polynomial bound
on $g$ gives a polynomial reduction. Membership in NP is
Proposition~\ref{thm:sat-np}. A 3-CNF is already a propositional formula,
so this same reduction establishes hardness of SAT. The intermediate
circuit compiler establishes hardness of circuit SAT.

Here ``3-CNF'' means clauses of width \emph{at most} three. If exactly
three literal occurrences are required, repeat a literal in a shorter
nonempty clause. If three distinct variables are required, replace
$(a\vee b)$ by $(a\vee b\vee z)\wedge(a\vee b\vee\neg z)$ with fresh
$z$, and replace a unit $a$ by the four clauses
$(a\vee z^{\epsilon}\vee u^{\delta})$ for
$\epsilon,\delta\in\{0,1\}$ with fresh $z,u$ and signed literals
$z^1=z,z^0=\neg z$. For an empty clause use all eight clauses on three
fresh variables, which cannot be simultaneously true. Tautological clauses
are deleted and repetitions simplified before these replacements. Direct
bit cases prove extension and restriction in each case. The size increase
is a fixed factor and hence remains polynomial.
\end{proof}

\begin{corollary}[Tautology is coNP-complete]\label{thm:taut-conp}
Relative to the same explicit verifier presentations for complements,
the tautology problem is coNP-complete under polynomial many-one reductions.
\end{corollary}
\begin{proof}
A nontautology has a polynomial-size falsifying valuation, so its complement
is in NP. If the complement of a problem reduces to satisfiability of
$\Phi_x$, then membership in the original problem is equivalent to
unsatisfiability of $\Phi_x$, hence to tautologicity of $\neg\Phi_x$.
Adding one negation is a polynomial transformation. The witness correctness
of the reduction was proved above, and the SAT/TAUT duality is
Theorem~\ref{thm:truth-decision}.
\end{proof}

This proof supplies no superpolynomial lower bound for SAT. NP-completeness
is a statement about polynomial reductions, fully compatible with the
primitive recursive decision algorithms already constructed.
